\documentclass[11pt]{article}
\usepackage[T1]{fontenc}
\usepackage{amsmath,amssymb,booktabs}
\usepackage[margin=1in]{geometry}
\usepackage[hidelinks]{hyperref}
\usepackage{tabularx}
\title{Proof Dossier: Weighted Frobenius Polarization}
\author{Project Koios Formal}
\date{}
\begin{document}
\maketitle
\section{Identity and status}
\begin{tabularx}{\textwidth}{@{}l>{\raggedright\arraybackslash}X@{}}
\toprule
Statement ID & \texttt{BG-HILBERT-THM-009} \\
Status & \texttt{proved} \\
Lean declaration & \path{ProjectKoios.BandGauge.Hilbert.rankTwoWeightedFrobeniusSq_linear_combination} \\
Module & \path{ProjectKoios/BandGauge/Hilbert/WeightedFrobenius.lean} \\
\bottomrule
\end{tabularx}

\section{Exact statement}
For real $a,b$ and finite complex rank-two fields $X,Y$,
\[
 P(aX+bY)=a^2P(X)+2ab\,C(X,Y)+b^2P(Y),
\]
where $P$ and $C$ use the same source coefficient
$\alpha_i/(2E_{\mathrm{ref}}^2)$.

\section{Source and proof}
This theorem \texttt{adapts} the Frobenius expansion used in Supplementary Section S1.1
of Ragasa2026A. Owner source: \path{ksdft2effmass}, immutable commit
\begin{center}\small\texttt{99a46cb59d0c6f27ada7d709f194996ef9619d9f}.\end{center}
The source file is
\path{docs/publications/papers/Ragasa2026a/v3/}\path{supplementary.tex}, lines 51--82.
The proof expands the squared complex modulus of each entry and then distributes the
finite weighted sum.

\section{Boundary and trusted base}
The identity is algebraic and does not instantiate scientific matrices. Lean and pinned
Mathlib report \texttt{propext}, \texttt{Classical.choice}, and \texttt{Quot.sound}.
No admission or scientific axiom is used. The proof and prose were AI-assisted and
kernel-checked. No independent expert review is recorded. Status: \texttt{proved};
working draft.
\end{document}
